\section*{Hoofdstuk \ref{ch:b2:microbial-metabolism} --- Microbieel metabolisme en biofilms}

\begin{solution}{exo:b2:microbial-metabolism:1}
Cyanobacterie: foto-litho-autotroof (licht, water, $\mathrm{CO_2}$).
\emph{Nitrosomonas}: chemo-litho-autotroof (oxidatie van ammonium,
$\mathrm{CO_2}$). \emph{E.~coli} op glucose: chemo-organo-heterotroof.
Paarse niet-zwavelbacterie op succinaat in het licht:
foto-organo-heterotroof. \emph{Thiobacillus} op sulfide in het donker:
chemo-litho-autotroof.
\end{solution}

\begin{solution}{exo:b2:microbial-metabolism:2}
Zuurstof (de bovenste millimeters), nitraat (net daaronder, waar de
zuurstof op is), mangaan- en ijzeroxiden (de overgang van bruin naar
grijs), sulfaat (de zwarte sulfidelaag, dik in mariene sedimenten),
$\mathrm{CO_2}$ voor de \omterm{def:b2:microbial-metabolism:anaerobic}{methanogenese} (het diepst, of net onder het
oppervlak in sulfaatarm zoet water).
\end{solution}

\begin{solution}{exo:b2:microbial-metabolism:3}
$5\,\mathrm{C_6H_{12}O_6} + 24\,\mathrm{NO_3^-} + 24\,\mathrm{H^+}
\to 30\,\mathrm{CO_2} + 12\,\mathrm{N_2} + 42\,\mathrm{H_2O}$: 4.8
nitraationen per glucose (elk nitraat neemt 5 elektronen op, elke
glucose staat er 24 af).
\end{solution}

\begin{solution}{exo:b2:microbial-metabolism:4}
Een aan een oppervlak vastzittende microbiële gemeenschap, ingebed in
een zelfgemaakte matrix van polysachariden, eiwitten en DNA. Stadia:
hechting, microkolonie en uitscheiding van matrix, rijping met torens en
kanalen, verspreiding. Tandplak en cariës; infecties van katheters en
implantaten (en van de longen bij taaislijmziekte); aangroei en corrosie
van leidingen en scheepsrompen --- of, nuttig, de oxidatiebedden van de
waterzuivering.
\end{solution}

\begin{solution}{exo:b2:microbial-metabolism:5}
Acht elektronen vallen van \qty{-0.42}{V} naar \qty{-0.24}{V}:
$\Delta G^{\circ\prime} = -8\times 96.5\times 0.18 =
\qty{-139}{kJ/mol}$ methaan (in tabellen \qty{-131}{kJ/mol}), d.w.z.\
\qty{35}{kJ} per $\mathrm{H_2}$: hoogstens 0.7 ATP per waterstof, in de
praktijk ongeveer een halve. De energie per substraatmolecuul is zo
klein dat methanogenen langzaam groeien (\omterm{def:b2:unicellular-diversity:division}{generatietijden} van uren tot
dagen) en bijna al hun substraat in gas omzetten in plaats van in
cellen.
\end{solution}

\begin{solution}{exo:b2:microbial-metabolism:6}
$\mu = 1.2\,S/(5 + S)$: \qty{0.20}{h^{-1}} ($T = \ln 2/\mu =
\qty{3.5}{h}$), \qty{0.60}{h^{-1}} (\qty{1.2}{h}), \qty{1.09}{h^{-1}}
(\qty{0.64}{h}). Negentig procent: $S = 9K_s = \qty{45}{mg/L}$.
\end{solution}

\begin{solution}{exo:b2:microbial-metabolism:7}
$S^* = 0.05\times 0.4/(0.8 - 0.4) = \qty{0.05}{g/L}$; $X^* =
0.4\times(5 - 0.05) = \qty{1.98}{g/L}$; productie $DX^* = 0.4\times
1.98 = \qty{0.79}{g/L/h}$; uitspoeling $D_c = 0.8\times 5/5.05 =
\qty{0.79}{h^{-1}}$.
\end{solution}

\begin{solution}{exo:b2:microbial-metabolism:8}
Vastgelegd: $0.25\times 275 = \qty{69}{kJ}$, d.w.z.\ 1.4 ATP per
ammonium; $12/1.4 = 8.7$, ongeveer negen ammoniumionen per koolstof --- en
meermaals meer zodra de kosten meetellen van het maken van NADH door
elektronen vanaf nitriet bergop te drijven. Een heterotroof krijgt zijn
koolstof al gereduceerd en georganiseerd, en dertig ATP per glucose; een
nitrificeerder moet zowel zijn energie als zijn reductievermogen bij een
zwakke donor kopen en elke koolstof uit $\mathrm{CO_2}$ opbouwen, zodat
hij slechts een piepkleine fractie van het omgezette substraat in cellen
omzet.
\end{solution}

\begin{solution}{exo:b2:microbial-metabolism:9}
$L_p = \sqrt{2\times 1.2\times 10^{-9}\times 0.2/0.5} =
\sqrt{9.6\times 10^{-10}} = \qty{31}{\micro m}$. Verdubbelen van $c_0$
vermenigvuldigt $L_p$ met $\sqrt 2$: \qty{44}{\micro m}; verviervoudigen
van $q$ halveert haar: \qty{15}{\micro m}.
\end{solution}

\begin{solution}{exo:b2:microbial-metabolism:10}
Van boven naar beneden: water met algen en cyanobacteriën (oxygene
fotosynthese); een roestkleurige band van ijzeroxideerders waar
$\mathrm{Fe^{2+}}$ zuurstof ontmoet; een witte sluier van kleurloze
zwaveloxideerders (chemolithotrofen) waar sulfide zuurstof ontmoet; een
paarse band van paarse zwavelbacteriën (\omterm{def:b2:microbial-metabolism:anoxygenic}{anoxygene fotosynthese} op
sulfide); een groene band van groene zwavelbacteriën (hetzelfde, met
meer tolerantie voor sulfide en minder licht); zwarte modder van gisters
en sulfaatreduceerders. (a) Paarse zwavelbacteriën hebben zowel licht
van boven als sulfide van onderen nodig, dus zitten ze op het grensvlak
waar de twee gradiënten elkaar kruisen. (b) Sulfaatreduceerders maken
$\mathrm{H_2S}$, dat ijzer neerslaat als zwart FeS. (c) Gesloten voor
materie: zwavel, koolstof en stikstof wisselen tussen geoxideerde en
gereduceerde vormen zonder de kolom te verlaten; open voor energie:
licht komt binnen, warmte gaat naar buiten, en zonder licht stopt elke
kringloop.
\end{solution}

\begin{solution}{exo:b2:microbial-metabolism:11}
Productie $pn$, verlies $kA$: $A^* = pn/k$. $A_c = 10^{-8}\times
6.02\times 10^{23} = 6.0\times 10^{15}$ moleculen per liter; $n_c =
kA_c/p = 5\times 10^{-4}\times 6.0\times 10^{15}/5 = 6\times 10^{11}$
per liter, $6\times 10^{8}$ per milliliter. De dichte cultuur
($10^{9}$) zit boven de drempel, zeewater ($10^{5}$) vier ordes van
grootte eronder. In een \omterm{def:b2:microbial-metabolism:biofilm}{biofilm} zitten de cellen op $10^{11}$ per
milliliter matrix, en de matrix vertraagt het ontsnappen van het signaal
(wat $k$ in feite verlaagt), zodat een microkolonie van enkele duizenden
cellen in een nanoliter het quorum bereikt, terwijl dezelfde cellen,
verspreid over een liter, dat nooit zouden doen.
\end{solution}

\begin{solution}{exo:b2:microbial-metabolism:12}
Stikstofbinding (alleen prokaryoten hebben nitrogenase): zonder haar zou
de gebonden stikstof van de biosfeer door \omterm{def:b2:microbial-metabolism:anaerobic}{denitrificatie} weglekken en
zou de productie binnen decennia ineenstorten. Nitrificatie en
\omterm{def:b2:microbial-metabolism:anaerobic}{denitrificatie} (alleen prokaryoten): zonder deze processen zou ammonium
zich ophopen en nitraat verdwijnen, en zou stikstof nooit naar de lucht
terugkeren. \omterm{def:b2:microbial-metabolism:anaerobic}{Methanogenese} en \omterm{def:b2:microbial-metabolism:anaerobic}{anaerobe ademhaling} (alleen prokaryoten):
zonder deze processen zou organisch materiaal zich in zuurstofloze
sedimenten, pensen en moerassen onafgebroken ophopen en zou koolstof de
kringloop verlaten; evenzo was oxygene fotosynthese een bacteriële
uitvinding, en zonder haar zou er geen zuurstof zijn om te ademen. De
eukaryoten zijn een late en chemisch beperkte toevoeging aan een
prokaryote planeet.
\end{solution}

\begin{solution}{pb:b2:microbial-metabolism:1}
\textbf{1.} $-24\times 96.5\times (0.82 + 0.43) = \qty{-2900}{kJ/mol}$.
\textbf{2.} Nitraat: $-24\times 96.5\times 1.17 = \qty{-2700}{kJ/mol}$;
sulfaat: $-24\times 96.5\times 0.21 = \qty{-490}{kJ/mol}$.
\textbf{3.} Zuurstof levert meer op, dus verdringen aeroben de
denitrificeerders bij de strijd om het organisch materiaal zolang er
zuurstof is, en zuurstof onderdrukt de nitraatreductasen; \omterm{def:b2:microbial-metabolism:anaerobic}{sulfaatreductie}
vereist zuurstofloosheid \emph{en} de afwezigheid van nitraat, dus haar
geur betekent dat de tank te weinig wordt belucht en het nitraat is
uitgeput --- en $\mathrm{H_2S}$ is giftig en tast het beton aan.
\textbf{4.} $0.4\times 2900/50 = 23$ ATP; $0.4\times 2700/50 = 22$;
$0.4\times 490/50 = 4$.
\textbf{5.} Ongeveer zes keer meer ($23/4$).
\textbf{6.} Twee ATP per glucose betekent een opbrengst van ongeveer
\num{0.1}, zodat \qty{90}{\%} van de energie van het substraat in de
producten blijft (zuren, alcoholen, $\mathrm{H_2}$), die de methanogenen
met een even kleine winst in methaan omzetten; weinig energie, weinig
biomassa, en de koolstof vertrekt als gas.
\textbf{7.} $10^4\times 0.3 = \qty{3000}{kg}$ CZV per dag.
\textbf{8.} Biomassa $0.4\times 3000 = \qty{1200}{kg}$ CZV; de
resterende \qty{1800}{kg} CZV wordt geoxideerd en vraagt \qty{1800}{kg}
$\mathrm{O_2}$.
\textbf{9.} $1800/2 = \qty{900}{kWh}$ per dag.
\textbf{10.} $10^4\times 0.04 = \qty{400}{kg}$ stikstof;
$4.57\times 400 = \qty{1830}{kg}$ $\mathrm{O_2}$.
\textbf{11.} In de stationaire toestand is $\mu = 1/\theta \le \mu_{\max}$:
uitspoeling onder $\theta = 1/\mu_{\max} = \qty{2}{d}$.
\textbf{12.} $D = \qty{0.1}{d^{-1}}$: $S^* = 1\times 0.1/(0.5 - 0.1)
= \qty{0.25}{g/m^3}$.
\textbf{13.} $S^* = 0.1/(0.25 - 0.1) = \qty{0.67}{g/m^3}$, en de
uitspoelgrens stijgt tot \qty{4}{d}; de operator moet de slibleeftijd
verlengen (minder slib spuien): bij $\theta = \qty{20}{d}$ is $S^* =
0.05/0.2 = \qty{0.25}{g/m^3}$ opnieuw.
\textbf{14.} $2.86\times 400 = \qty{1144}{kg}$ CZV per dag.
\textbf{15.} Ammonium in het effluent $0.25\times 10^4 = \qty{2.5}{kg/d}$,
\qty{0.6}{\%} van de vracht; de rest, \qty{397}{kg} stikstof,
vertrekt als $\mathrm{N_2}$: \qty{99}{\%}.
\textbf{16.} Bespaard: \qty{1144}{kg} $\mathrm{O_2}$. Totale behoefte:
$(1800 - 1144) + 1830 = \qty{2490}{kg}$ $\mathrm{O_2}$ per dag.
\textbf{17.} $1200\times 0.6\times 0.35 = \qty{252}{m^3}$ methaan
per dag.
\textbf{18.} $252\times 36 = \qty{9070}{MJ} = \qty{2520}{kWh}$;
elektriciteit $0.35\times 2520 = \qty{880}{kWh}$ per dag.
\textbf{19.} Beluchting $2490/2 = \qty{1245}{kWh}$ per dag: het methaan
dekt daarvan $880/1245 = \qty{71}{\%}$.
\textbf{20.} $D_e\,c'' = q$, dus $c = c_0 + ax + qx^2/2D_e$; met
$c(L_p) = 0$ en $c'(L_p) = 0$: $a = -qL_p/D_e$ en $c =
(q/2D_e)(L_p - x)^2$, met $L_p^2 = 2D_ec_0/q$.
\textbf{21.} $L_p = \sqrt{2\times 1.5\times 10^{-9}\times 0.2/0.3} =
\sqrt{2\times 10^{-9}} = \qty{45}{\micro m}$.
\textbf{22.} $45/300 = \qty{15}{\%}$ aeroob. Daaronder reduceren
denitrificeerders het nitraat dat vanuit de nitrificerende oppervlaktelaag
omlaag diffundeert tot $\mathrm{N_2}$, met het organisch materiaal dat
naar binnen diffundeert.
\textbf{23.} Binnen één film schept de zuurstofgradiënt een aerobe laag
(nitrificatie) boven een zuurstofloze laag (\omterm{def:b2:microbial-metabolism:anaerobic}{denitrificatie}), enkele
tientallen micrometers van elkaar; de beluchte tank is overal zuurstofrijk,
dus vraagt \omterm{def:b2:microbial-metabolism:anaerobic}{denitrificatie} daar een aparte zuurstofloze tank.
\textbf{24.} Chloor reageert met de matrixpolymeren en wordt erdoor
verbruikt voordat het de diepe cellen bereikt; de diepe cellen zijn
uitgehongerd, groeien traag of zijn in rust en zijn veel minder gevoelig;
alleen de oppervlaktelaag sterft af en de film groeit van onderen weer aan.
\textbf{25.} Zuurstofbehoefte \qty{2490}{kg} per dag; methaan
\qty{252}{m^3} per dag; het methaan levert \qty{71}{\%} van de
beluchtingselektriciteit (\qty{880}{kWh} van \qty{1245}{kWh}).
\end{solution}
